\section{Related Work}
\label{sec:related}

Deep learning lineage studies broadly infer model relationships from either \emph{model weights} or \emph{behavior}. These methods have been driven by different applications including provenance and IP \citep{horwitz2025mother, zeng2024huref, nikolic2026model}, model poisoning and misalignment \citep{cloud2025subliminal}, trait inheritance \citep{cloud2025subliminal, agents2026same}, and data auditing \citep{zhu2025independence}. 

\textbf{Model weights:} Several works trace weight changes from fine-tuning or reinforcement learning from human feedback (RLHF) back to parent models. MoTHer \citep{horwitz2025mother} reconstructs directed model trees from weight similarity via a directed minimum spanning tree and evaluates against controlled, known model hierarchies; \citet{yu2025neural} identify parent--child and grandparent--child lineage with both a learning-free fine-tuning approximation and a learned detector ($>99\%$ accuracy). Beyond direct lineage, spectral weight ``fingerprints'' identify common ancestry \citep{ghostspec2025}, and Spectral DeTuning \citep{horwitz2024spectral} reconstructs base weights from shared LoRA variants. Unlike MoTHer, which reconstructs relationships among observed parent and descendant models, we observe only leaf models and infer a phylogeny with latent ancestral branching points. This setting is relevant to LLM family reconstruction, where intermediate or parent models may not be available.

REEF \citep{zhang2025reef} compares internal representations using centered kernel alignment (CKA; \citealp{kornblith2019similarity}) to identify model derivation, while TensorGuard \citep{tensorguard2025} uses gradient-response fingerprints for pairwise model similarity and family classification.

Other works examine \emph{why} lineage is detectable in models. Models fine-tuned from a shared initialization stay within a single low-error basin and are robust to SGD noise \citep{neyshabur2020,wortsman2022model,frankle2020linear}, so pairwise weight distance can approximate genealogical distance. Additionally, task vectors \citep{ilharco2023editing} are near-orthogonal across tasks, so distinct lineages contribute independent directions. LoRA \citep{hu2022lora} and intrinsic-dimensionality results \citep{aghajanyan2021intrinsic} demonstrate that meaningful adaptation can occur in low-dimensional parameter subspaces. Together, these results motivate the possibility that lineage information remains detectable after fine-tuning.

\textbf{Behavior:} Behavioral approaches infer model relationships using only API access or model outputs, without requiring access to model weights \citep{yax2025phylolm,wu2026llmdna}. PhyloLM \citep{yax2025phylolm} builds LLM phylogenies from output similarity and predicts benchmark performance, while LLM DNA \citep{wu2026llmdna} defines a functional ``DNA'' representation carrying lineage information. Both methods estimate phylogenies over hundreds of models. Other output-only methods test pairwise provenance from next-token distributions \citep{nikolic2026model} or output independence \citep{zhu2025independence}. A key advantage of behavioral methods is that they can compare models across architectures and sizes, whereas weight-based comparisons generally require compatible parameterizations. Conversely, weight-based lineage estimates do not depend on the choice of prompts or test tasks.

\textbf{Our Approach:} We compare weight- and behavior-based approaches, determine where phylogenetic signal can be detected, and evaluate the estimated phylogenies against known true trees. Our findings are consistent with prior work showing that inherited behavioral signals can be preserved under sufficient training signal \citep{cloud2025subliminal}, as well as phylogenetic results showing that the underlying tree can be correctly estimated when certain conditions on the estimated distances are satisfied \citep{atteson1999performance}.  As Table~\ref{tab:related} shows, ours is the only work that combines weight-based estimation, latent ancestors (which allow for unobserved intermediate models), controlled LLM fine-tuning trees, a link between weights and behavior, and a label-free tree-likeness score that serves as an empirical measure of confidence in the estimated phylogeny.

\begin{table}[t]
\caption{We compare with the closest prior methods. Controlled trees are families the authors fine-tuned. Empirical confidence is a label-free score that is tested against estimation accuracy.}
\label{tab:related}
\begin{center}
\small
\setlength{\tabcolsep}{4pt}
\begin{tabular}{lcccccc}
\toprule
Method & Access & \shortstack{Latent\\ancestors} & \shortstack{Controlled\\trees} & \shortstack{Weight--\\behavior} & \shortstack{Empirical\\confidence} & \shortstack{Cross-\\arch.} \\
\midrule
MoTHer \citep{horwitz2025mother}
  & weights  & \ding{55}  & vision     & \ding{55}  & \ding{55}  & \ding{55} \\
Neural Phylogeny \citep{yu2025neural}
  & weights  & \ding{55}  & vision     & \ding{55}  & \ding{55}  & \ding{55} \\
PhyloLM \citep{yax2025phylolm}
  & behavior & \checkmark & \ding{55}  & \ding{55}  & \ding{55}  & \checkmark \\
LLM DNA \citep{wu2026llmdna}
  & behavior & \checkmark & \ding{55}  & \ding{55}  & \ding{55}  & \checkmark \\
\textbf{Ours}
  & \textbf{weights} & \checkmark & \textbf{LLM} & \checkmark & \checkmark & \ding{55} \\
\bottomrule
\end{tabular}
\end{center}
\end{table}
